\documentclass[11pt]{article}
\usepackage{amsmath,amssymb,booktabs,graphicx}
\usepackage[margin=1in]{geometry}

\title{Initial Experiments: Extending the MVP Simulator}
\author{William Dennis}
\date{June 2026}

\begin{document}
\maketitle

Extends the MVP simulator (\texttt{docs/mvp/mvp\_specification.tex}) by enriching
the state space, action space, and reward function, moving toward the full
MDP defined in \texttt{docs/design/initial\_design.tex}.

\section{State Space Extension}
\label{sec:state}

References \texttt{docs/design/initial\_design.tex \S3} and Issue \#124.


\subsection{Enriched StateView}

% To be filled: archetype, burden, goal_progress, steps, previous_action


\subsection{Contextual Bandit Conditioning}

% To be filled: context_feature config, conditional action selection


\subsection{Open Questions}

% To be filled


\section{Action Space Extension}
\label{sec:action}

References \texttt{docs/design/initial\_design.tex \S3} and Issue \#126.


\subsection{ActionConfig}

% To be filled: reward_penalty, burden_penalty, 6-action set


\subsection{Transition Model Updates}

% To be filled: 6x6x2 transition matrix, backward compatibility


\subsection{Agent Compatibility}

% To be filled: verification of variable-length action handling


\subsection{Open Questions}

% To be filled


\section{Reward Function Extension}
\label{sec:reward}

References \texttt{docs/design/initial\_design.tex \S3} and Issue \#132.
\emph{Implementation: PR 1c (feat/1c-reward) --- merged.}


\subsection{RewardWeightsConfig}

An optional \texttt{reward\_weights} block controls the delayed reward
component. When absent the reward is purely per-state (MVP backward-compatible
mode):
\[
r(s, a, t) = m[t \bmod S] \cdot r(s')
\]

When present, the schema accepts these fields:

\begin{center}
\begin{tabular}{lll}
\toprule
Field & Type & Default \\
\midrule
\texttt{mode} & \texttt{Literal["multi\_timescale"]} & \texttt{"multi\_timescale"} \\
\texttt{delayed\_reward\_interval} & \texttt{int $\ge$ 1} & 21 \\
\texttt{delayed\_reward\_value} & \texttt{float $\ge$ 0} & 10.0 \\
\texttt{delayed\_reward\_scale} & \texttt{float $\mid$ None} & None \\
\texttt{delayed\_reward\_threshold} & \texttt{float $\in$ [0,1] $\mid$ None} & None \\
\bottomrule
\end{tabular}
\end{center}

All numeric fields are validated at load time via Pydantic \texttt{Field}
constraints. The mode is constrained to \texttt{Literal["multi\_timescale"]};
when additional modes are added later the union can be expanded.

\subsection{Flat Bonus (Option A)}

When \texttt{delayed\_reward\_scale} is \texttt{None} a flat bonus fires at
every interval boundary:
\[
r_{\text{bonus}} =
\begin{cases}
\texttt{value} & \text{if } t > 0 \text{ and } t \equiv 0 \pmod{N} \\
0 & \text{otherwise}
\end{cases}
\]
where $N = \texttt{delayed\_reward\_interval}$ and $t$ is the zero-indexed
cumulative step count (\texttt{state.global\_step}).

\subsection{Scaled Bonus (Option B)}

When \texttt{delayed\_reward\_scale} is set the bonus depends on the fraction
of steps in the interval where the simulated user was in the active state:
\[
\text{rate} = \frac{\text{active\_count}}{N}
\qquad
r_{\text{bonus}} =
\begin{cases}
\texttt{scale} \times \text{rate} & \text{if } \text{rate} \ge \texttt{threshold} \\
0 & \text{otherwise}
\end{cases}
\]

The active count is tracked in \texttt{CompoundReward.\_active\_count},
incremented on every \texttt{reward()} call where
\texttt{state.activity == "active"} and reset to zero at each interval
boundary.

\subsection{Episode Boundaries}

The \texttt{RewardHandler} ABC now defines a \texttt{reset()} method that
clears episode-local state (\texttt{\_active\_count}).
\texttt{Environment.reset()} calls \texttt{self.\_reward.reset()} so that
reusing a reward instance across episodes does not leak state.

\subsection{Results}

The 1c experiment config is at
\texttt{docs/initial\_experiments/configs/1c\_reward\_only.yaml}.
Mean total rewards across 50 seeds (seeds 0--49) for each agent and
reward mode:

\begin{center}
\small
\begin{tabular}{lrrrrr}
\toprule
& \multicolumn{5}{c}{Mean total reward (50 seeds)} \\
\cmidrule(lr){2-6}
Agent & No bonus & Flat bonus & Scaled bonus & Flat $-$ base & Scaled $-$ base \\
\midrule
Thompson Sampling & 160.1 & 369.5 & 163.0 & 209.4 & 2.8 \\
Random            & 136.3 & 346.3 & 138.0 & 210.0 & 1.8 \\
$\epsilon$-greedy & 157.4 & 350.5 & 160.4 & 193.1 & 3.0 \\
UCB               & 155.4 & 364.8 & 159.0 & 209.4 & 3.6 \\
Decaying $\epsilon$ & 156.4 & 360.3 & 158.7 & 203.9 & 2.2 \\
\bottomrule
\end{tabular}
\end{center}

The flat bonus (interval=21, value=10.0) adds roughly 200--210 in
total reward across all agents, consistent with the expected
$21$--$22$ firings $\times 10.0 = 210$ bonus over a 450-step episode.
The scaled bonus (scale=5.0, threshold=0.6) adds only 2--4 on average:
no agent sustains a 60\% active rate within a 21-step interval given
the hand-tuned transition probabilities. Thompson Sampling achieves
the highest base reward (160.1) and benefits most from the flat bonus
(369.5 total).

\subsection{Open Questions}

\begin{itemize}
\item The \texttt{delayed\_reward\_interval} and
\texttt{delayed\_reward\_value} parameters are placeholders. Calibration
against real-world user data would inform clinically meaningful values.
\item The scaled bonus threshold (default~0.6) and scale (used as 5.0 in this run; schema default is \texttt{None}) were
chosen arbitrarily. A grid search could identify regions where the delayed
reward shapes agent behaviour without overwhelming the per-step signal.
\item The flat vs.\ scaled choice is a config toggle today. A future extension
could blend both: a fixed bonus plus a scaled bonus for users who meet the
threshold.
\end{itemize}


\section{Combined Results}

Results from the pairwise and full-integration experiments will appear here
once all three extension PRs (1a, 1b, 1c) have been merged and the integration
PRs (4--7) are complete. The config matrix of seven experiment configurations
is defined in \texttt{docs/initial\_experiments/plans/state\_action\_reward\_epic.md}.
Currently available:

\begin{itemize}
\item \texttt{docs/initial\_experiments/configs/1c\_reward\_only.yaml}
(PR~1c, 2~actions, 2~states, multi-timescale reward) --- see \S\ref{sec:reward}.
\end{itemize}

\end{document}
